\documentclass[12pt,a4paper,french]{article}

\usepackage[utf8]{inputenc}
\usepackage[T1]{fontenc}
\usepackage{babel}
\usepackage{geometry}
\usepackage{xcolor}
\usepackage{hyperref}
\usepackage{graphicx}
\usepackage{listings}
\usepackage{titlesec}
\usepackage{tcolorbox}
\usepackage{enumitem}
\usepackage{booktabs}
\usepackage{longtable}
\usepackage{tabularx}
\usepackage{amsmath}
\usepackage{setspace}
\usepackage{parskip}

\geometry{
    top=2.5cm, bottom=2.5cm,
    left=2.5cm, right=2.5cm
}

\definecolor{primaryblue}{RGB}{31,78,121}
\definecolor{lightblue}{RGB}{214,228,240}
\definecolor{accentblue}{RGB}{70,130,180}
\definecolor{darkgray}{RGB}{64,64,64}
\definecolor{lightgray}{RGB}{245,245,245}
\definecolor{successgreen}{RGB}{0,128,0}
\definecolor{warningorange}{RGB}{204,102,0}
\definecolor{errored}{RGB}{180,30,30}
\definecolor{codebg}{RGB}{248,248,248}

\hypersetup{
    colorlinks=true,
    linkcolor=primaryblue,
    urlcolor=accentblue,
    pdftitle={Note Méthodologique — Calcul des Salaires},
    pdfauthor={Data Engineering Team},
}

\titleformat{\section}[block]
    {\normalfont\Large\bfseries\color{primaryblue}}
    {\thesection}{1em}{}

\titleformat{\subsection}[block]
    {\normalfont\large\bfseries\color{accentblue}}
    {\thesubsection}{1em}{}

\titleformat{\subsubsection}[block]
    {\normalfont\normalsize\bfseries\color{darkgray}}
    {\thesubsubsection}{1em}{}

\tcbset{
    common/.style={
        enhanced, breakable,
        fonttitle=\bfseries\color{primaryblue},
        left=6pt, right=6pt, top=4pt, bottom=4pt
    }
}

\newtcolorbox{infobox}[1][]{
    common,
    colback=lightblue!30, colframe=primaryblue,
    title={#1}
}

\newtcolorbox{warningbox}[1][]{
    common,
    colback=orange!10, colframe=warningorange,
    title={#1}
}

\newtcolorbox{infoboxsimple}{
    enhanced, breakable,
    colback=lightblue!30, colframe=primaryblue,
    left=6pt, right=6pt, top=4pt, bottom=4pt
}

\lstset{
    backgroundcolor=\color{codebg},
    basicstyle=\small\ttfamily,
    breaklines=true,
    captionpos=b,
    commentstyle=\color{successgreen},
    keywordstyle=\color{primaryblue}\bfseries,
    stringstyle=\color{errored},
    numberstyle=\tiny\color{darkgray},
    numbers=left,
    numbersep=8pt,
    frame=single,
    framesep=4pt,
    rulecolor=\color{darkgray!40},
    showstringspaces=false,
    tabsize=4,
    language=Python,
}

\newcommand{\colname}[1]{\texttt{#1}}
\newcommand{\code}[1]{\texttt{#1}}

\setlength{\parskip}{6pt}
\setlength{\parindent}{0pt}
\onehalfspacing

\begin{document}

% ─── Page de titre ───────────────────────────────────────────
\begin{titlepage}
    \centering
    \vspace*{2cm}
    {\huge\bfseries\color{primaryblue}
        Note Méthodologique\par
        \vspace{0.5cm}
        Calcul des Salaires\par
        \vspace{0.3cm}
        Panel Administratif 2015--2025\par
    \par}
    \vspace{1.5cm}
    \textcolor{accentblue}{\rule{10cm}{2pt}}
    \vspace{1cm}

    {\large\bfseries Version 1.0\par}
    \vspace{0.5cm}
    {\large Guide méthodologique complet sur le calcul du salaire brut,\par
           les arbitrages de population et les imputations.\par}
    \vspace{2cm}

    \begin{tcolorbox}[colback=lightblue!40, colframe=primaryblue, width=12cm]
        \centering
        \textbf{Prérequis :} Comprendre la structure F1/F2 des fichiers source\par
        \textbf{Public :} Data engineers, statisticiens, responsables métier\par
    \end{tcolorbox}

    \vfill
    {\small\color{darkgray} Document généré le \today\par}
\end{titlepage}

\newpage

% ─── Table des matières ──────────────────────────────────────
\tableofcontents
\newpage

% ═════════════════════════════════════════════════════════════
\section{Contexte et objectifs}
% ═════════════════════════════════════════════════════════════

\subsection{Problématique}

Les fichiers de solde mensuels de la fonction publique contiennent deux feuilles :

\begin{itemize}
    \item \textbf{Feuille 1 (F1)} : données administratives de l'agent
          (matricule, nom, organisme, grade, montant brut agrégé, etc.)
    \item \textbf{Feuille 2 (F2)} : décomposition détaillée par code de
          rémunération (rubriques 124, 228, 271, 421, 425, etc.)
\end{itemize}

Cette structure soulève plusieurs questions méthodologiques :

\begin{enumerate}
    \item Quelle feuille utiliser comme source du salaire ?
    \item Comment gérer les individus présents dans une feuille mais pas l'autre ?
    \item Comment traiter les individus sans composante de salaire dans F2 ?
    \item Comment imputer les grades manquants ?
\end{enumerate}

\subsection{Objectifs de cette note}

Ce document décrit :

\begin{itemize}
    \item Les choix méthodologiques retenus pour le calcul du salaire
    \item Les arbitrages faits sur la population d'étude
    \item Les règles d'imputation des grades
    \item Les règles d'imputation des salaires
    \item Les contrôles qualité appliqués
\end{itemize}

% ═════════════════════════════════════════════════════════════
\section{Structure des fichiers source}
% ═════════════════════════════════════════════════════════════

\subsection{Description des feuilles}

\begin{tabularx}{\textwidth}{lX}
    \toprule
    \textbf{Feuille} & \textbf{Contenu} \\
    \midrule
    \textbf{Feuille 1 (F1)} &
        Données agent : matricule, nom, date naissance, sexe,
        situation administrative, organisme, lieu d'affectation,
        emploi, grade, \textbf{montant\_brut} (agrégé),
        \textbf{montant\_net} (agrégé), etc. \\
    \textbf{Feuille 2 (F2)} &
        Décomposition par code numérique : chaque ligne représente un
        individu $\times$ code poste avec un montant.
        La clé de jointure est \texttt{MATRICULE||CODE\_ORGANISME}. \\
    \bottomrule
\end{tabularx}

\subsection{Codes rubriques F2}

Les codes numériques de la Feuille 2 correspondent à des types de rémunération :

\begin{tabularx}{\textwidth}{lX}
    \toprule
    \textbf{Code} & \textbf{Libellé (ANSTAT)} \\
    \midrule
    \code{124\_F2} & Solde de Base \\
    \code{228\_F2} & Indemnité de transport / complément \\
    \code{271\_F2} & Prime de rendement ou équivalent \\
    \code{421\_F2} & Majoration ou indemnité spécifique \\
    \code{425\_F2} & Autre élément de rémunération \\
    \bottomrule
\end{tabularx}

\begin{warningbox}[Codes obligatoires]
    Les codes \code{124\_F2} et \code{421\_F2} sont considérés comme
    indispensables. Un individu sans ces deux codes ne peut pas être inclus
    dans l'analyse salariale.
\end{warningbox}

% ═════════════════════════════════════════════════════════════
\section{Arbitrages méthodologiques}
% ═════════════════════════════════════════════════════════════

\subsection{Choix 1 : INNER JOIN plutôt que LEFT JOIN}

\subsubsection{Problème}

原始ement, le pipeline utilisait un \texttt{LEFT JOIN} :
F1 conservait tous les individus, avec des valeurs nulles pour F2
lorsque la jointure échouait.

Cela posait deux problèmes :

\begin{enumerate}
    \item \textbf{Population incohérente} : des individus présents uniquement
          en F1 (sans données F2) étaient inclus avec un salaire calculable
          uniquement via le champ \texttt{montant\_brut} de F1.
    \item \textbf{Deux sources de vérité} : \texttt{montant\_brut} (F1)
          pouvait différer de la somme des composantes F2.
\end{enumerate}

\subsubsection{Décision}

Depuis la version 2.1 du pipeline :

\begin{itemize}
    \item \textbf{INNER JOIN} : seuls les individus présents dans les deux
          feuilles sont conservés
    \item \textbf{Filtrage F2} : dans F2, on ne garde que les individus
          ayant \code{124\_F2} ET \code{421\_F2} non nuls
\end{itemize}

\begin{infoboxsimple}
    \textbf{Raison} : Le salaire brut doit être calculé à partir des
    composantes F2, pas du champ agrégé F1. Un individu sans ces composantes
    ne peut pas avoir un salaire fiable.
\end{infoboxsimple}

\subsection{Choix 2 : Suppression du fallback}

\subsubsection{Problème}

Le code précédent contenait un \texttt{fallback} :

\begin{lstlisting}[language=Python]
if composantes_presentes:
    # somme normale
else:
    # fallback vers montant_brut de F1
    base = base.with_columns(
        pl.col("montant_brut").cast(pl.Float64).alias("SALAIRE_BRUT")
    )
\end{lstlisting}

Ce fallback masquait les problèmes de données et introduisait une
\textbf{incohérence méthodologique} : certains salaires venaient de F2,
d'autres de F1.

\subsubsection{Décision}

Depuis la version 2.1 :

\begin{lstlisting}[language=Python]
if composantes_presentes:
    # somme des composantes F2
else:
    raise RuntimeError(
        "Aucune composante de salaire trouvee. "
        "Verifiez que l'etape 02 a bien ete executee."
    )
\end{lstlisting}

\begin{warningbox}[Conséquence]
    Ce changement nécessite de \textbf{ré-exécuter le pipeline complet}
    (étapes 02 à 08) pour recalculer les données avec la nouvelle logique.
\end{warningbox}

% ═════════════════════════════════════════════════════════════
\section{Calcul du salaire brut}
% ═════════════════════════════════════════════════════════════

\subsection{Rationnel méthodologique}

La méthodologie de calcul des salaires vise à établir des moyennes comparables
entre secteurs en s'affranchissant des spécificités de chaque grille salariale.

\begin{infobox}[Principe de comparabilité]
    Pour réduire l'effet des particularités de chaque grille salariale sur la
    moyenne, le calcul du salaire moyen des fonctionnaires se fait en tenant
    compte uniquement des codes postes de salaires que \textbf{tous les
    fonctionnaires ont en commun}. Les indemnités spécifiques à l'emploi et/ou
    à la fonction du fonctionnaire sont retirées.
\end{infobox}

\begin{warningbox}[Exclusion du code poste 135]
    Le code poste \texttt{135} ne doit pas être intégré dans la détermination
    des moyennes spécifiques par secteur. Ce code représente des indemnités
    spécifiques qui ne sont pas communes à l'ensemble des fonctionnaires.
\end{warningbox}

\subsection{Codes postes communs retenus}

Les codes postes communs à tous les fonctionnaires et retenus pour le calcul
du salaire brut sont :

\begin{center}
\begin{tabular}{cll}
    \toprule
    \textbf{Code} & \textbf{Libellé (ANSTAT)} & \textbf{Rôle} \\
    \midrule
    \code{124\_F2} & Solde de Base & Composante principale du traitement \\
    \code{228\_F2} & Indemnité de transport / complément & Composante commune \\
    \code{271\_F2} & Prime de rendement ou équivalent & Composante commune \\
    \code{421\_F2} & Majoration ou indemnité spécifique & Composante commune \\
    \code{425\_F2} & Autre élément de rémunération & Composante commune \\
    \bottomrule
\end{tabular}
\end{center}

\subsection{Formule de calcul}

\[
\text{SALAIRE\_BRUT} = 124\_F2 + 228\_F2 + 271\_F2 + 421\_F2 + 425\_F2
\]

\begin{infobox}[Justification de la formule]
    C'est la somme de ces codes postes qui est considérée comme salaire brut
    en vue de déterminer la moyenne générale par grade des fonctionnaires.
    Cette approche garantit que les comparaisons entre secteurs et entre grades
    ne sont pas biaisées par des indemnités spécifiques à certaines fonctions.
\end{infobox}

\subsection{Logigramme de calcul}

\begin{enumerate}
    \item \textbf{Étape 02 — Jointure et filtrage}
          \begin{itemize}
              \item Lecture de F1 et F2 depuis le fichier Excel mensuel
              \item Filtrage F2 : garder uniquement les individus avec
                    \code{124\_F2} ET \code{421\_F2} non nuls
              \item INNER JOIN sur \texttt{MATRICULE||CODE\_ORGANISME}
          \end{itemize}

    \item \textbf{Étape 07 — Calcul}
          \begin{itemize}
              \item Cast des composantes en \texttt{Float64}
              \item Somme horizontale des composantes disponibles
              \item Si aucune composante : erreur explicite (plus de fallback)
          \end{itemize}
\end{enumerate}

\subsection{Indicateurs produits}

À partir de \colname{SALAIRE\_BRUT}, le pipeline produit :

\begin{tabularx}{\textwidth}{lX}
    \toprule
    \textbf{Indicateur} & \textbf{Définition} \\
    \midrule
    \colname{REVENU\_TOTAL} & Somme des \colname{SALAIRE\_BRUT} par agent $\times$ période \\
    \colname{SALAIRE\_BRUT\_MOYEN} & Moyenne par grade $\times$ mois \\
    \colname{SALAIRE\_BRUT\_MEDIAN} & Médiane par grade $\times$ mois \\
    \bottomrule
\end{tabularx}

% ═════════════════════════════════════════════════════════════
\section{Imputation des grades}
% ═════════════════════════════════════════════════════════════

\subsection{Problématique}

Le champ \colname{statut\_fonctionnaire} (F1) contient des libellés textuels
de la forme \texttt{"Catégorie A3"} ou \texttt{"Classe B2"}. Certains
enregistrements ont un grade manquant (\texttt{NF}) ou invalide (grade
incomplet comme \texttt{A} ou grade inexistant comme \texttt{B7}).

\textbf{Contexte historique :} Avant 2021, les grades étaient mal renseignés
dans les fichiers de solde. Par exemple, un individu pouvait avoir \texttt{B7}
en 2018 alors que son grade réel était \texttt{B3}. Depuis 2021, le
renseignment s'est amélioré.

\subsection{Extraction du grade}

Le grade est extrait par expression régulière :

\begin{lstlisting}[language=Python]
def extraire_grade(statut):
    if not statut:
        return "NF"
    match = re.search(r"Cat[eé]gorie\s+([ABCD][0-9]*)", str(statut), re.I)
    if match:
        return match.group(1).upper().strip()  # ex: "A3"
    return "NF"
\end{lstlisting}

\subsection{Classification des grades}

\begin{center}
\begin{tabular}{lll}
    \toprule
    \textbf{Type} & \textbf{Grades} & \textbf{Traitement} \\
    \midrule
    Valides & A1--A7, B1--B6, C1--C5, D1--D3 & Conservés tels quels \\
    Invalides & B7, C7, D7, A, B, C, D, NF & Imputation forward-fill \\
    Non imputables & NF sans historique & Exclus des indicateurs \\
    \bottomrule
\end{tabular}
\end{center}

\subsection{Méthode : Forward-Fill par individu (étape 03b)}

L'imputation utilise une approche \textbf{forward-fill individuel} :

\begin{enumerate}
    \item \textbf{Principe}
          \begin{itemize}
              \item Pour un grade invalide en année $t$, on regarde si le même
                    individu a un grade valide dans une année future $t' > t$
              \item L'individu est identifié par \colname{cle\_unique} +
                    \colname{emploi}
              \item Si l'emploi n'a pas changé, on reporte le grade futur
                    vers le passé
          \end{itemize}

    \item \textbf{Exemple}
          \begin{itemize}
              \item Individu \texttt{XYZ} en 2018 : grade = \texttt{B7} (invalide)
              \item Individu \texttt{XYZ} en 2022 : grade = \texttt{B3} (valide),
                    même emploi
              \item $\Rightarrow$ Le grade \texttt{B3} est reporté sur 2018
          \end{itemize}

    \item \textbf{Construction du lookup}
          \begin{itemize}
              \item On parcourt toute la période (2015--2025)
              \item Pour chaque individu (\colname{cle\_unique} +
                    \colname{emploi}), on garde le grade valide de l'année
                    la plus récente
              \item Ce grade futur est utilisé pour imputer les années antérieures
          \end{itemize}

    \item \textbf{NF\_NON\_IMPUTABLE}
          \begin{itemize}
              \item Si l'individu n'a aucun grade valide dans les années futures,
                    le grade reste invalide
              \item Ces individus sont exclus des calculs d'indicateurs
          \end{itemize}
\end{enumerate}

\begin{infobox}[Pourquoi le forward-fill individuel ?]
    Plutôt que d'utiliser le grade modal d'un emploi (qui peut concerner
    plusieurs personnes), le forward-fill garde la trace de chaque individu.
    Si une personne est \texttt{B3} en 2022, c'est son propre grade qui est
    reporté en arrière, pas celui d'un autre fonctionnaire occupant le même
    emploi.
\end{infobox}

\subsection{Source du grade}

\begin{tabularx}{\textwidth}{lX}
    \toprule
    \textbf{Valeur} & \textbf{Signification} \\
    \midrule
    \texttt{OBSERVE} & Grade extrait directement de \colname{statut\_fonctionnaire},
                       valide \\
    \texttt{IMPUTE} & Grade reporté depuis une année future du même individu \\
    \texttt{NF\_NON\_IMPUTABLE} & Aucun grade valide trouvé dans l'historique
                       de l'individu \\
    \bottomrule
\end{tabularx}

% ═════════════════════════════════════════════════════════════
\section{Imputation des salaires}
% ═════════════════════════════════════════════════════════════

\subsection{Situation}

Certains individus ont un \colname{SALAIRE\_BRUT} nul ou manquant.
Cela peut survenir si :

\begin{itemize}
    \item Les composantes F2 sont à zéro pour certains mois
    \item L'individu est en \texttt{demi\_solde} (mi-temps administratif)
\end{itemize}

\subsection{Traitement}

Le paramètre \colname{INCLURE\_ZEROS} contrôle le comportement :

\begin{tabularx}{\textwidth}{lX}
    \toprule
    \textbf{Valeur} & \textbf{Comportement} \\
    \midrule
    \texttt{True} & Les salaires à zéro ou NA sont inclus dans les calculs \\
    \texttt{False} & Seuls les salaires strictement > 0 sont inclus \\
    \bottomrule
\end{tabularx}

\begin{warningbox}[Impact sur les indicateurs]
    Inclure les zéros réduit la moyenne et la médiane. Cette option est
    pertinente pour analyser la \textbf{masse salariale totale}.
    L'exclusion des zéros est plus adaptée pour analyser les
    \textbf{salaires effectivement versés}.
\end{warningbox}

% ═════════════════════════════════════════════════════════════
\section{Winsorisation}
% ═════════════════════════════════════════════════════════════

\subsection{Problématique}

Les fichiers de solde peuvent contenir des valeurs extrêmes (erreurs de saisie,
indemnités exceptionnelles, etc.) qui faussent les indicateurs de tendance
centrale.

\subsection{Méthode}

La winsorisation replace les valeurs en dehors d'un intervalle par les
bornes de cet intervalle :

\[
x_{wins} = \begin{cases}
    P_1 & \text{si } x < P_1 \\
    x & \text{si } P_1 \leq x \leq P_{99} \\
    P_{99} & \text{si } x > P_{99}
\end{cases}
\]

où $P_1$ et $P_{99}$ sont les percentiles 1 et 99.

\subsection{Paramètres}

\begin{itemize}
    \item \colname{SEUIL\_WINS} = 0.01 (par défaut) $\rightarrow$ écrêtage à P1/P99
    \item \colname{MIN\_OBS\_WINS} = 100 (par défaut) : effectif minimum requis
          pour calculer les percentiles par groupe CITP
\end{itemize}

\subsection{Niveaux de winsorisation}

\begin{enumerate}
    \item \textbf{Par \texttt{mois × grade × CITP}} : si effectif $\geq$ 100
    \item \textbf{Par \texttt{mois × grade}} (repli) : si effectif $<$ 100
\end{enumerate}

\begin{infobox}[Winsorisation en mémoire]
    La winsorisation est appliquée uniquement en mémoire pour le calcul
    des indicateurs. Les données Gold ne sont pas modifiées.
\end{infobox}

% ═════════════════════════════════════════════════════════════
\section{Filtrage par situation administrative}
% ═════════════════════════════════════════════════════════════

\subsection{Situations normalisées}

\begin{tabularx}{\textwidth}{lX}
    \toprule
    \textbf{Code} & \textbf{Libellés reconnus} \\
    \midrule
    \texttt{en\_activite} & EN ACTIVITE, ACTIVITE, EN ACTIVITÉ, ACTIVITÉ \\
    \texttt{regul\_indemnites} & REGUL. INDEMNITES, REGUL INDEMNITES,
                                  REGULARISATION INDEMNITES \\
    \texttt{demi\_solde} & DEMI-SOLDE, DEMI SOLDE, DEMISOLDE, 1/2 SOLDE \\
    \texttt{autre} & Toute autre valeur (incluant null et vide) \\
    \bottomrule
\end{tabularx}

\subsection{Traitement}

Les situations \texttt{autre} sont \textbf{supprimées} à l'étape 03.
Elles ne participent pas au calcul des indicateurs.

% ═════════════════════════════════════════════════════════════
\section{Contrôles qualité}
% ═════════════════════════════════════════════════════════════

\subsection{Contrôles sur les grades}

\begin{itemize}
    \item Distribution OBSERVE / IMPUTE / NF\_NON\_IMPUTABLE affichée
    \item Alerte si taux de NF\_NON\_IMPUTABLE > 5\%
\end{itemize}

\subsection{Contrôles sur les montants}

\begin{itemize}
    \item Comptage des lignes où \texttt{montant\_net > montant\_brut}
          (incohérence : le net ne peut pas dépasser le brut)
    \item Distribution des salaires > 0
\end{itemize}

\subsection{Contrôles de matching ANSTAT}

\begin{itemize}
    \item Taux de matching par code (ORGANISME, GRADE, EMPLOI, etc.)
    \item Logs en temps réel pour les codes surveillés
\end{itemize}

% ═════════════════════════════════════════════════════════════
\section{Récapitulatif des étapes de calcul}
% ═════════════════════════════════════════════════════════════

\subsection{Flux complet}

\begin{enumerate}
    \item \textbf{02 — Staging vers Bronze}
          \begin{itemize}
              \item Lecture F1 et F2
              \item Filtrage F2 : 124\_F2 ET 421\_F2 non nuls
              \item INNER JOIN F1 $\leftrightarrow$ F2
          \end{itemize}

    \item \textbf{03 — Bronze vers Silver}
          \begin{itemize}
              \item Mapping colonnes vers noms standardisés
              \item Normalisation situations administratives
              \item Filtrage : suppression des \texttt{autre}
              \item Enrichissement codes ANSTAT
              \item Cast types numériques
          \end{itemize}

    \item \textbf{05 — Silver vers Gold}
          \begin{itemize}
              \item Agrégation masse\_salariale.parquet
              \item Agrégation effectifs.parquet
          \end{itemize}

    \item \textbf{06 — Compilation panel}
          \begin{itemize}
              \item Construction panel individu $\times$ période
              \item Déduplication matricule $\times$ mois
          \end{itemize}

    \item \textbf{06b — Imputation grades}
          \begin{itemize}
              \item Extraction grade depuis statut\_fonctionnaire
              \item Cascade P1 $\rightarrow$ P2 $\rightarrow$ P3
              \item Marquage source (OBSERVE/IMPUTE/NF)
          \end{itemize}

    \item \textbf{07 — Calcul indicateurs}
          \begin{itemize}
              \item Calcul SALAIRE\_BRUT = somme composantes F2
              \item Winsorisation P1/P99
              \item Agrégation REVENU\_TOTAL par agent
              \item Calcul indicateurs (moy, méd, min, max, P25, P75)
          \end{itemize}

    \item \textbf{08 — Export Excel}
          \begin{itemize}
              \item Génération fichier multi-feuilles
              \item Création sommaire navigable
          \end{itemize}
\end{enumerate}

% ═════════════════════════════════════════════════════════════
\section{Glossaire}
% ═════════════════════════════════════════════════════════════

\begin{tabularx}{\textwidth}{lX}
    \toprule
    \textbf{Terme} & \textbf{Définition} \\
    \midrule
    \textbf{Bronze} & Données brutes, tout en String, telles qu'ingérées \\
    \textbf{Silver} & Données enrichies, typées, avec codes ANSTAT \\
    \textbf{Gold} & Données agrégées ou panelisées, prêtes pour l'analyse \\
    \textbf{CITP} & Classification Internationale Type des Professions \\
    \textbf{INNER JOIN} & Jointure conservant uniquement les lignes communes \\
    \textbf{Winsorisation} & Remplacement des valeurs extrêmes par les percentiles \\
    \bottomrule
\end{tabularx}

\vfill

\begin{center}
    \textcolor{accentblue}{\rule{10cm}{1pt}}

    \smallskip
    \textbf{Document :} Note Méthodologique — Calcul des Salaires

    \textbf{Version :} 1.0

    \textbf{Date :} \today

    \smallskip
    \textit{Pour toute question, contacter l'équipe Data Engineering.}
\end{center}

\end{document}
